\documentclass[11pt]{article}
% ======================================================================
%  examstyle.tex — EPFL-style exam layout for the weekly Time Series exams
%  Shared by every Exam_Week_NN(.tex / _Solutions.tex). Compiles with pdflatex.
% ======================================================================
\usepackage[utf8]{inputenc}
\usepackage[T1]{fontenc}
\usepackage[english]{babel}
\usepackage[a4paper,margin=2.4cm]{geometry}
\usepackage{amsmath,amssymb,amsthm,mathtools}
\usepackage{bm}
\usepackage{enumitem}
\usepackage{array}

\setlength{\parindent}{0pt}
\setlength{\parskip}{0.5em}
\emergencystretch=3em
\allowdisplaybreaks[1]

% ---------- math macros (same names as the rest of the project) ----------
\newcommand{\E}{\mathbb{E}}
\DeclareMathOperator{\Var}{Var}
\DeclareMathOperator{\Cov}{Cov}
\DeclareMathOperator{\Corr}{Corr}
\DeclareMathOperator*{\argmin}{arg\,min}
\DeclareMathOperator{\MSE}{MSE}
\DeclareMathOperator{\trace}{tr}
\newcommand{\Reals}{\mathbb{R}}
\newcommand{\Ints}{\mathbb{Z}}
\newcommand{\Comp}{\mathbb{C}}
\newcommand{\Nat}{\mathbb{N}}
\newcommand{\Prob}{\mathbb{P}}
\newcommand{\Tr}{^{\mathsf{T}}}
\newcommand{\conj}[1]{\overline{#1}}
\newcommand{\abs}[1]{\left\lvert #1 \right\rvert}
\newcommand{\norm}[1]{\left\lVert #1 \right\rVert}
\newcommand{\ind}{\mathbf{1}}
\newcommand{\eps}{\varepsilon}
\newcommand{\diff}{\nabla}
\newcommand{\acvs}{\gamma}
\newcommand{\acf}{\rho}
\newcommand{\pacf}{\alpha}
\newcommand{\sdf}{\mathcal{S}}
\newcommand{\Bsh}{B}
\newcommand{\iid}{\overset{\text{iid}}{\sim}}

\setlist[enumerate]{leftmargin=2em,topsep=2pt,itemsep=4pt}
\setlist[itemize]{leftmargin=1.6em,topsep=2pt,itemsep=2pt}

\newcounter{exo}

% ---------- exam cover header :  \examheader{exam title}{duration} ----------
\newcommand{\examheader}[2]{%
  \thispagestyle{empty}
  \begin{center}
    {\Large\bfseries Time Series}\\[3pt]
    Spring Semester 2026, EPFL\\
    \'Ecole Polytechnique F\'ed\'erale de Lausanne\\
    Prof.\ Dr.\ Sofia C.\ Olhede\\[7pt]
    \hrule height 0.8pt
    \vspace{9pt}
    {\large\bfseries Practice Exam --- #1}\\[4pt]
    {\normalsize Duration: #2 \quad$\cdot$\quad No calculator \quad$\cdot$\quad Answer on paper}\\
    \vspace{8pt}
    \hrule height 0.8pt
  \end{center}
  \vspace{14pt}
  \begin{tabular}{@{}p{8.2cm}@{}p{8cm}@{}}
    Name:\enspace\hrulefill & Surname:\enspace\hrulefill
  \end{tabular}\\[14pt]
  SCIPER (Student Card Number):\enspace\hrulefill
  \vspace{16pt}
}

% ---------- points table (fixed: 4 problems) ----------
\newcommand{\pointstable}{%
  \begin{center}
  \renewcommand{\arraystretch}{1.5}
  \begin{tabular}{|p{5cm}|p{4cm}|}
    \hline
    \textbf{Exercise} & \textbf{Points}\\\hline
    1 & \\\hline
    2 & \\\hline
    3 & \\\hline
    4 & \\\hline
    \textbf{Total} & \\\hline
  \end{tabular}
  \end{center}
}

% ---------- problem heading (exam: one problem per page) ----------
\newcommand{\problem}[1]{%
  \clearpage
  \stepcounter{exo}%
  \begin{center}\textbf{\large Problem \theexo\ (#1 points)}\end{center}
  \vspace{4pt}
}
% ---------- answer space: fill the statement page + one blank answer page ----------
% Put \answerpage at the END of each problem so every exercise gets its own page(s).
\newcommand{\answerpage}{\par\vfill\newpage\null\newpage}

% ---------- solutions document ----------
\newcommand{\solheader}[1]{%
  \begin{center}
    {\Large\bfseries Time Series --- Solutions}\\[3pt]
    {\large Practice Exam --- #1}\\[2pt]
    EPFL, MATH-342 \quad$\cdot$\quad Prof.\ S.\ C.\ Olhede\\[5pt]
    \hrule height 0.8pt
  \end{center}
  \vspace{10pt}
}
\newcommand{\solproblem}[1]{%
  \bigskip
  \stepcounter{exo}%
  {\large\bfseries Problem \theexo\ (#1 points)}\par\nobreak\vspace{2pt}
}
% Rappel de l'énoncé avant la solution (dans les corrigés)
\newcommand{\statementstart}{\par\vspace{1pt}{\bfseries\itshape Statement.}\par\nobreak\begingroup\small\itshape}
\newcommand{\statementend}{\par\endgroup\vspace{5pt}\hrule\vspace{6pt}{\bfseries Solution.}\par\nobreak\vspace{2pt}}

% --- local macros not in examstyle ---
\newcommand{\filt}{L}
\begin{document}

\examheader{Week 13: Linear Time-Invariant (LTI) Filters}{60 minutes}
\pointstable
\begin{center}\itshape Justify every answer. Throughout, $\eps_t$ is a mean-zero white noise of variance $\sigma^2$.\end{center}

% ---------------------------------------------------------------
\problem{5}
\textbf{(LTI axioms, impulse response and transfer function.)} A digital filter $\filt$ maps an input sequence to an output sequence, $\{Y_t\}=\filt[\{X_t\}]$; we write a single output entry as $Y_u=\filt[\{X_t\}]_u$. Recall that $\filt$ is \emph{linear time-invariant} (LTI) if it is linear and commutes with the backshift, $\filt\Bsh=\Bsh\filt$. The impulse response is $h_m=\filt[\{\delta_{t,-m}\}]_0$ and, for $h\in\ell^1$, the transfer function is $H(f)=\sum_{m\in\Ints}h_m e^{-2\pi i fm}$.

\textbf{(a)} Let $\filt_1,\filt_2$ be LTI filters and $\alpha\in\Comp$. Prove, directly from the two axioms, that the linear combination $\filt=\alpha\filt_1+\filt_2$ is again LTI (check linearity and time invariance separately). \hfill(2 pts)

\textbf{(b)} Consider the concrete one-sided filter
\[
\filt[\{X_t\}]_u=X_u-\tfrac12\,X_{u-1},\qquad u\in\Ints .
\]
Identify its impulse response $\{h_m\}$, write $\filt$ as a polynomial $\Phi(\Bsh)$ in the backshift, and compute its transfer function $H(f)$ together with the squared gain $\abs{H(f)}^2$ in real form (no complex exponentials). \hfill(2 pts)

\textbf{(c)} At which frequency $f\in[0,\tfrac12]$ is the squared gain $\abs{H(f)}^2$ of part (b) \emph{largest}, and at which is it \emph{smallest}? Give both extreme values of $\abs{H(f)}^2$, and state in one sentence whether this filter is low-pass or high-pass. \hfill(1 pt)
\answerpage

% ---------------------------------------------------------------
\problem{5}
\textbf{(Differencing, smoothing, and the output spectrum.)} Let $\{X_t\}$ be a zero-mean stationary process with $\acvs^{(X)}\in\ell^1$ and spectral density $\sdf_X(f)$. Recall the output-spectrum theorem: if $\{Y_t\}=\filt[\{X_t\}]$ for an LTI filter with transfer function $H$, then $\sdf_Y(f)=\abs{H(f)}^2\sdf_X(f)$.

\textbf{(a)} Let $\{Y_t\}=(I-\Bsh)[\{X_t\}]$, i.e.\ $Y_t=X_t-X_{t-1}$. Explain briefly why $\{Y_t\}$ is stationary, then prove that
\[
\sdf_Y(f)=4\sin^2(\pi f)\,\sdf_X(f).
\]
(Show the algebraic step $\abs{1-e^{-2\pi i f}}^2=4\sin^2(\pi f)$.) \hfill(2 pts)

\textbf{(b)} Now cascade the difference with the centred three-point average
\[
\filt_2[\{V_t\}]_u=\tfrac13\bigl(V_{u-1}+V_u+V_{u+1}\bigr),
\]
forming $\{W_t\}=\filt_2\bigl[(I-\Bsh)[\{X_t\}]\bigr]$. Using the fact that the transfer function of a cascade is the \emph{product} of the two transfer functions, express $\sdf_W(f)$ in terms of $\sdf_X(f)$, and name every frequency in $[0,\tfrac12]$ at which $\sdf_W$ is forced to vanish. \hfill(2 pts)

\textbf{(c)} Specialise to white-noise input, $\sdf_X(f)\equiv\sigma^2$, and evaluate $\sdf_Y(\tfrac14)$ for the differenced series of part (a). \hfill(1 pt)
\answerpage

% ---------------------------------------------------------------
\problem{5}
\textbf{(ARMA spectral density and time reversibility of AR$(p)$.)} For a stationary $\mathrm{ARMA}(p,q)$ process $\Phi(\Bsh)X_t=\Theta(\Bsh)\eps_t$ the spectral density is
\[
\sdf_X(f)=\sigma^2\,\frac{\abs{\Theta(e^{-2\pi i f})}^2}{\abs{\Phi(e^{-2\pi i f})}^2},
\qquad
\Phi(z)=1-\sum_{j=1}^p\phi_j z^j,\quad \Theta(z)=1-\sum_{j=1}^q\theta_j z^j .
\]

\textbf{(a)} For the AR$(1)$ process $X_t=\phi X_{t-1}+\eps_t$ with $\abs\phi<1$, use the polynomial-filter lemma (the filter $\Phi(\Bsh)$ has transfer function $\Phi(e^{-2\pi i f})$) to derive
\[
\sdf_X(f)=\frac{\sigma^2}{1+\phi^2-2\phi\cos(2\pi f)} ,
\]
and state, with reason, where $\sdf_X$ peaks when $\phi>0$ and when $\phi<0$. \hfill(2 pts)

\textbf{(b)} Let $\{X_t\}$ be a Gaussian AR$(p)$ process, $X_t=\sum_{j=1}^p\phi_j X_{t-j}+\eps_t$, and define the time-reversed series $Y_t=X_{-t}$. Show that $\acvs^{(Y)}_\tau=\acvs^{(X)}_\tau$ for all $\tau$, hence $\sdf_Y(f)=\sdf_X(f)$. \hfill(1.5 pts)

\textbf{(c)} Put $\phi_0=-1$ and $\tilde\nu_t=\sum_{j=0}^p\phi_j Y_{t-j}=\Phi(\Bsh)[\{Y_t\}]_t$. Using parts (a)--(b), the polynomial-filter lemma and the output-spectrum theorem, prove that $\sdf_{\tilde\nu}(f)\equiv\sigma^2$, and conclude that $\{Y_t\}$ is an AR$(p)$ with the \emph{same} coefficients $\phi_1,\dots,\phi_p$ driven by white noise of variance $\sigma^2$. \hfill(1.5 pts)
\answerpage

% ---------------------------------------------------------------
\problem{5}
\textbf{(Revision of Week 12: ARCH$(1)$ --- moments, white noise, and the ACF of the squares.)} Let $\{r_t\}$ be an $\mathrm{ARCH}(1)$ process,
\[
r_t=\sigma_t\eps_t,\qquad \sigma_t^2=\alpha_0+\alpha_1 r_{t-1}^2,
\]
with $\alpha_0>0$, $0\le\alpha_1<1$, and $\{\eps_t\}$ i.i.d.\ standard Gaussian (mean $0$, variance $1$) independent of $\mathcal F_{t-1}$.

\textbf{(a)} Using that $\sigma_t$ is $\mathcal F_{t-1}$-measurable and iterated expectation, show $\E[r_t]=0$, and (assuming weak stationarity) that
\[
\Var(r_t)=\frac{\alpha_0}{1-\alpha_1}.
\]
\hfill(1.5 pts)

\textbf{(b)} Prove that $\{r_t\}$ is serially uncorrelated, i.e.\ $\Cov(r_{t+\tau},r_t)=0$ for all $\tau\neq0$, by conditioning on $\mathcal F_{t+\tau-1}$. (This is what makes ARCH a \emph{white noise}.) \hfill(1.5 pts)

\textbf{(c)} Assume in addition $0<\alpha_1<1/\sqrt3$ (so a finite fourth moment exists). Recall $\Corr(r_{t+\tau}^2,r_t^2)=\alpha_1^{\abs\tau}$ and the kurtosis $\kappa=\dfrac{3(1-\alpha_1^2)}{1-3\alpha_1^2}$. Take $\alpha_1=\tfrac12$ and report the numerical value of $\Corr(r_{t+1}^2,r_t^2)$, of $\Corr(r_{t+2}^2,r_t^2)$, and of the kurtosis $\kappa$; state in one sentence why $\kappa>3$ means fat tails. \hfill(2 pts)
\answerpage

\end{document}
